% SGA 3, Expose IV, English reconstruction.
% Scope: opening through Lemma 1.6 and its proof, stopping before Lemma 1.7.
% Authority: Polo--Gille Exp4-13oct24.pdf, component-local pages 1--3
% (combined-reader pages 189--191; printed pages 160--163).
% Jacob Reinhold's CC BY 4.0 Markdown lineage was used only as a phrasing
% comparison. The Polo--Gille PDF controls prose, formulas, numbering, and notes.

\setcounter{footnote}{0}
\SGARefTarget{sga3:iv:chapter:iv}\chapter*{第 IV 篇\\拓扑与层}
\addcontentsline{toc}{chapter}{第 IV 篇. 拓扑与层}
\markboth{第 IV 篇. 拓扑与层}%
  {第 IV 篇. 拓扑与层}

\begingroup
\renewcommand{\thefootnote}{\fnsymbol{footnote}}
\begin{center}
  \textit{M. Demazure 著}\footnote{本文展开了 A.~Grothendieck 两次讲课的实质内容,
  并在若干原先略而未论或仅简略提及的重要之处作了补充.}
\end{center}
\endgroup

\setcounter{footnote}{-1}

% Source page: Exp4 component-local p. 1; combined-reader p. 189; printed p. 160.
本篇旨在使读者熟悉拓扑与层的语言的基本内容 (不涉及上同调).
这种语言对于构造商等问题尤其方便.
\footnote{编者注: 2024 年 10 月 13 日版本.}

前三节建立构造商所需的语言. 第四节是本篇的核心, 介绍层论,
主要着眼于它在商的问题中的应用; 第五节将这一理论应用于群中的商的构造
以及主齐性丛. 最后一节更具体地研究概形范畴, 定义其上的若干有用拓扑.

读者可以参阅 \SGARefLink{sga3:iv:bib:as}{[AS]},
\SGARefLink{sga3:iv:bib:ma}{[MA]}, \SGARefLink{sga3:iv:bib:d}{[D]}
和 \SGARefLink{sga3:iv:bib:sga:4}{SGA~4}; 特别地,
关于拓扑在下降理论中的应用, 可参阅 \SGARefLink{sga3:iv:bib:d}{[D]};
关于宇宙的问题, 可参阅 \SGARefLink{sga3:iv:bib:sga:4}{SGA~4}
(本篇对这方面的处理尤其欠妥).

\SGARefTarget{sga3:iv:section:1}\section*{1. 普遍有效满态射}
\addcontentsline{toc}{section}{1. 普遍有效满态射}

在本篇余下部分中, 固定一个范畴 \(\mathscr C\).

\medskip
\noindent\SGARefTarget{sga3:iv:definition:1:1}\textbf{定义 1.1.}
\textit{若对每个对象 \(X\), 相应的映射}
\[
  X(S)=\operatorname{Hom}(S,X)
  \longrightarrow
  X(T)=\operatorname{Hom}(T,X)
\]
\textit{都是单射, 则称态射 \(u:T\to S\) 为\emph{满态射 (epimorphism)}.}
\footnote{编者注: 即 \(u\) 可右消去.}
\textit{若对每个态射 \(S'\to S\), 纤维积 \(T'=T\times_S S'\)
存在, 且 \(u':T'\to S'\) 为满态射, 则称它为
\emph{普遍满态射 (universal epimorphism)}.}

% Source page: Exp4 component-local p. 2; combined-reader p. 190; printed p. 162.
\medskip
\noindent\SGARefTarget{sga3:iv:definition:1:2}\textbf{定义 1.2.}
\textit{集合映射构成的图}
\[
  A\xrightarrow{u}B
  \overset{v_1}{\underset{v_2}{\rightrightarrows}}C
\]
\textit{称为\emph{正合的}, 是指 \(u\) 为单射, 且其像恰由满足
\(v_1(b)=v_2(b)\) 的元素 \(b\in B\) 组成.
\(\mathscr C\) 中同样形式的图称为\emph{正合的}, 是指对
\(\mathscr C\) 的每个对象 \(X\), 相应的集合图}
\[
  A(X)\longrightarrow B(X)\rightrightarrows C(X)
\]
\textit{都是正合的; 此时也称 \(u\) 使 \(A\) 成为箭头对
\((v_1,v_2)\) 的一个\emph{核}. 对偶地, \(\mathscr C\) 中的图}
\[
  C\overset{v_1}{\underset{v_2}{\rightrightarrows}}B
  \xrightarrow{u}A
\]
\textit{称为正合的, 是指它在反范畴 \(\mathscr C^{\circ}\)
中作为图是正合的; 即对 \(\mathscr C\) 的每个对象 \(X\),
相应的集合图}
\[
  X(A)\longrightarrow X(B)\rightrightarrows X(C)
\]
\textit{都是正合的.}\footnote{编者注: 特别地, 这蕴含 \(u\) 为满态射.}
\textit{此时也称 \(u\) 使 \(A\) 成为箭头对 \((v_1,v_2)\)
的一个\emph{余核}.}

\medskip
\noindent\SGARefTarget{sga3:iv:definition:1:3}\textbf{定义 1.3.}
\textit{若纤维平方 \(T\times_S T\) 存在, 且图}
\[
  T\times_S T
  \overset{\operatorname{pr}_1}{\underset{\operatorname{pr}_2}{\rightrightarrows}}
  T\xrightarrow{u}S
\]
\textit{正合, 则称态射 \(u:T\to S\) 为\emph{有效满态射
(effective epimorphism)}; 即 \(u\) 使 \(S\) 成为
\((\operatorname{pr}_1,\operatorname{pr}_2)\) 的余核.
若对每个态射 \(S'\to S\), 纤维积 \(T'=T\times_S S'\)
都存在, 且态射 \(u':T'\to S'\) 为有效满态射, 则称它为
\emph{普遍有效满态射 (universal effective epimorphism)}.}

显然有以下蕴含关系
\[
\begin{array}{ccc}
 \text{普遍有效满态射}
   &\Longrightarrow& \text{有效满态射}\\
 \Big\Downarrow && \Big\Downarrow\\
 \text{普遍满态射}
   &\Longrightarrow& \text{满态射},
\end{array}
\]
但一般没有其他蕴含关系成立.
\footnote{编者注: 例如, 若 \(\mathscr C=\mathbf{Sch}\) 为概形范畴,
容易看出每个普遍满态射都是满射. 令
\[
  T=\coprod_{p\ \text{素数}}\operatorname{Spec}(\mathbf F_p),
  \qquad S=\operatorname{Spec}(\mathbf Z).
\]
则态射 \(u:T\to S\) 是满态射, 却不是普遍满态射.
另一方面, \(T\times_S T\) 可等同于 \(T\), 因此两个投影
\(T\times_S T\to T\) 相同. 由于 \(\operatorname{id}_T\)
不能下降为态射 \(S\to T\), 可知 \(u\) 不是有效满态射.}

% Source page: Exp4 component-local p. 3; combined-reader p. 191; printed p. 163.
\medskip
\noindent\SGARefTarget{sga3:iv:definition:1:4:0}\textbf{定义 1.4.0.}
\footnote{编者注: 为便于后文引用, 添加了编号
\SGARefLink{sga3:iv:definition:1:4:0}{1.4.0}.}
我们``回顾'': 态射 \(u:T\to S\) 称为\emph{可平方的 (squarable)},
是指对每个态射 \(S'\to S\), 纤维积 \(T\times_S S'\) 都存在.

\medskip
\noindent\SGARefTarget{sga3:iv:lemma:1:4}\textbf{引理 1.4.}
\textit{考虑态射 \(U\xrightarrow{v}T\xrightarrow{u}S\). 则:}
\begin{enumerate}
  \item[(a)]\SGARefTarget{sga3:iv:lemma:1:4:item:a:01}
    \textit{若 \(u\) 和 \(v\) 是满态射, 则 \(uv\) 是满态射;
    若 \(uv\) 是满态射, 则 \(u\) 是满态射;}
  \item[(b)]\SGARefTarget{sga3:iv:lemma:1:4:item:b:01}
    \textit{若 \(u\) 和 \(v\) 是普遍满态射, 则 \(uv\) 是普遍满态射;
    若 \(uv\) 是普遍满态射且 \(u\) 可平方, 则 \(u\) 是普遍满态射.}
\end{enumerate}

引理~\SGARefLink{sga3:iv:lemma:1:4}{1.4} 由定义立即可得. 由此推出:

\medskip
\noindent\SGARefTarget{sga3:iv:corollary:1:5}\textbf{推论 1.5.}
\textit{设 \(u:X\to Y\) 和 \(u':X'\to Y'\) 是普遍满态射,
并假设 \(Y\times Y'\) 存在. 则 \(X\times X'\) 存在, 且}
\[
  u\times u':X\times X'\longrightarrow Y\times Y'
\]
\textit{是普遍满态射.}

还要记下以下定义和引理.

\medskip
\noindent\SGARefTarget{sga3:iv:definition:1:6:0}\textbf{定义 1.6.0.}
\footnote{编者注: 为便于后文引用, 添加了编号
\SGARefLink{sga3:iv:definition:1:6:0}{1.6.0}.}
\(\mathscr C\) 的对象 \(S\) 称为\emph{可平方的}, 是指它与
\(\mathscr C\) 的每个对象的乘积都存在. (若 \(\mathscr C\)
有终对象 \(e\), 这就等价于态射 \(S\to e\) 可平方;
参见 \SGARefLink{sga3:iv:definition:1:4:0}{1.4.0}.)

\medskip
\noindent\SGARefTarget{sga3:iv:lemma:1:6}\textbf{引理 1.6.}
\textit{设 \(u:X\to Y\) 是 \(\mathscr C_{/S}\) 中的态射.
它为满态射 (分别地, 普遍满态射、有效满态射或普遍有效满态射)
的充分条件是 \(\mathscr C\) 中的相应态射具有同一性质;
若假设 \(S\) 是 \(\mathscr C\) 的可平方对象, 则这也是必要条件.}

证明是直接的, 留给读者. \(S\) 可平方的假设用来将
从 \(\mathscr C_{/S}\) 的对象 \(Y\) 到 \(\mathscr C\) 的对象 \(Z\)
的 \(\mathscr C\)-态射, 解释为从 \(Y\) 到 \(Z\times S\) 的
\(\mathscr C_{/S}\)-态射.
% Translation statistics: source_lines=186; output_lines=181; inline_math=79; display_math=9; footnotes=7; pending_TODO=0.
% Review: complete assigned component; mathematical symbols, targets, references, tags, environments and footnotes checked; formula prose translated.
